\documentclass[a4paper,12pt]{extarticle}
\usepackage{geometry}
\usepackage[T1]{fontenc}
\usepackage[utf8]{inputenc}
\usepackage[english,russian]{babel}
\usepackage{amsmath}
\usepackage{amsthm}
\usepackage{amssymb}
\usepackage{fancyhdr}
\usepackage{setspace}
\usepackage{graphicx}
\usepackage{colortbl}
\usepackage{tikz}
\usepackage{pgf}
\usepackage{subcaption}
\usepackage{listings}
\usepackage{indentfirst}
\usepackage[
backend=biber,
style=numeric,
maxbibnames=99
]{biblatex}
\addbibresource{refs.bib}
\usepackage[colorlinks,citecolor=blue,linkcolor=blue,bookmarks=false,hypertexnames=true, urlcolor=blue]{hyperref} 
\usepackage{indentfirst}
\usepackage{mathtools}
\usepackage{booktabs}
\usepackage[flushleft]{threeparttable}
\usepackage{tablefootnote}

\usepackage{chngcntr} % нумерация графиков и таблиц по секциям
\counterwithin{table}{section}
\counterwithin{figure}{section}

\graphicspath{{graphics/}}%путь к рисункам

\makeatletter
% \renewcommand{\@biblabel}[1]{#1.} % Заменяем библиографию с квадратных скобок на точку:
\makeatother

\geometry{left=2.5cm}% левое поле
\geometry{right=1.0cm}% правое поле
\geometry{top=2.0cm}% верхнее поле
\geometry{bottom=2.0cm}% нижнее поле
\setlength{\parindent}{1.25cm}
\renewcommand{\baselinestretch}{1.5} % междустрочный интервал


\newcommand{\bibref}[3]{\hyperlink{#1}{#2 (#3)}} % biblabel, authors, year
\addto\captionsrussian{\def\refname{Список литературы (или источников)}} 

\renewcommand{\theenumi}{\arabic{enumi}}% Меняем везде перечисления на цифра.цифра
\renewcommand{\labelenumi}{\arabic{enumi}}% Меняем везде перечисления на цифра.цифра
\renewcommand{\theenumii}{.\arabic{enumii}}% Меняем везде перечисления на цифра.цифра
\renewcommand{\labelenumii}{\arabic{enumi}.\arabic{enumii}.}% Меняем везде перечисления на цифра.цифра
\renewcommand{\theenumiii}{.\arabic{enumiii}}% Меняем везде перечисления на цифра.цифра
\renewcommand{\labelenumiii}{\arabic{enumi}.\arabic{enumii}.\arabic{enumiii}.}% Меняем везде перечисления на цифра.цифра

\begin{document}
\begin{titlepage}
  \newpage

  {\setstretch{1.0}
  \begin{center}
    \universityname\\[0.6cm]
    \facultyname\\
    \programname
  \end{center}
  }

  \vspace{6em}

  \begin{center}
    {\bfseries Отчёт об исследовательском проекте на тему:}\\[0.5cm]
    {\bfseries \courseworktitle}\\[0.5cm]
    (промежуточный, этап 1)
  \end{center}

  \vspace{2em}

  {\bfseries Выполнил студент: \vspace{2mm}}

  {\setstretch{1.1}
  \begin{tabular}{l@{\hskip 1.5cm}l}
    группы \studentgroup, \studentcourse & \studentname
  \end{tabular}}

  \vspace{1em}
  {\bfseries Принял руководитель проекта: \vspace{2mm}}

  {\setstretch{1.1}
  \begin{tabular}{l}
    \supervisorname\\
    \supervisortitle
  \end{tabular}}

  \vfill

  \begin{center}
    \reportcity\ \reportyear
  \end{center}
\end{titlepage}
% это титульный лист - выберите подходящий вам из имеющихся в проекте вариантов (kr - курсовая работа у 3 курса, vkr - выпускная квалификационная работа у 4 курса)
\newpage
\setcounter{page}{2}

{
	\hypersetup{linkcolor=black}
	\tableofcontents
}

\newpage

\newpage
\section*{Аннотация}   % this is how to use russian
Работа посвящена исследованию методов семантической сегментации и классификации заболеваний растений с использованием foundation‑моделей компьютерного зрения для улучшения качества обнаружения патологий в реальных аграрных условиях.

Основное внимание уделено применению модели DINOv3 как базового визуального энкодера для выделения пораженных областей при ограниченной разметке (few‑shot learning) и анализу двух ключевых наборов данных: сегментационного PlantSeg с пиксельными масками и классификационного PlantVillage с метками классов.


\addcontentsline{toc}{section}{Аннотация}

\section*{Ключевые слова}
Глубинное обучение, классификация картинок, семантическая сегментация, foundation models, трансформер для картинок, DINOv3, few-shot learning, доменная адаптация, детекция болезней растений. 
\pagebreak

\section{Введение} 
\subsection{Описание предметной области и задачи}

Современное растениеводство сталкивается с рядом критически важных задач, среди которых автоматизированное выявление заболеваний растений и мониторинг их состояния занимают одно из центральных мест. Заболевания и физиологические стрессы негативно влияют на \textit{урожайность}, \textit{качество продукции} и \textit{экономическую эффективность} сельскохозяйственных предприятий. Традиционные методы диагностики, основанные на визуальном осмотре и вручную собранных данных, являются трудоемкими, субъективными и слабо масштабируемыми, поэтому возникает большая потребность в автоматизации этого процесса.

Семантическая сегментация представляет собой задачу классификации каждого пикселя изображения и его разделения на семантически значимые области. В контексте сельского хозяйства она позволяет \textit{локализовать и выделять }пораженные участки растений, что является необходимым шагом как для точечной диагностики, так и для дальнейшего автоматизированного принятия решений. В отличие от классических методов классификации изображений, которые дают только глобальные метки, семантическая сегментация предоставляет \textit{локальную, пиксель-уровневую информацию} о патологических изменениях, значительно повышая точность и информативность анализа.

\subsection{Использование foundation модели и адаптация}

Одним из современных подходов к построению универсальных визуальных представлений является использование самообучающихся моделей \textit{(self-supervised learning)}, которые способны извлекать высококачественные признаки из визуальных данных без необходимости в обширной аннотированной разметке. В частности, foundation модель \textbf{DINOv3} ~\cite{dinov3} (self DIstillation with NO labels v3) - это крупномасштабная визуальная модель на основе архитектуры \textit{ViT}, обученная методом самообучения с самодистилляцией на огромном количестве изображений \textit{(1.689 млрд. веб-картинок)}, что позволяет ей генерировать богатые плотные (dense) представления которые потом можно использовать во множестве различных downstream задач, в том числе сегментации.

Одной из распространенных проблем машинного обучения в растениеводстве является адаптация моделей к новым доменам с ограниченным количеством размеченных примеров. \textit{Few-shot} адаптация - подход, при котором модели обучаются эффективно обобщать на новые классы или условия с минимальным количеством размеченных примеров - становится особенно актуальной для сегментации заболеваний растений, где вручную размеченные данные часто дороги и труднодоступны.

\section{Обзор литературы}

\subsection{DINOv3 ~\cite{dinov3}}
Серия моделей DINO - прорывная в сфере компьютерного зрения, и последняя ее версия \textbf{DINOv3} кроме увеличения обучающих данных внесла новую технику \textit{Gram Anchoring}: во время обучения вводится дополнительный loss, заставляющий матрицу Грама \textit{cosine similarity} эмбеддингов патчей у модели ученика (в контексте \textit{self distillation}) быть схожей с матрицей учителя. Это позволяет модели лучше стабилизировать локальные фичи, продолжая иметь возможность их улучшать.

Так как качество локальных патчей крайне важно при задаче семантической сегментации, при выборе используемой модели в качестве backbone взгляд упал именно на нее.

\subsection{PlantVillage ~\cite{plantvillage}}

\textbf{PlantVillage} является одним из наиболее часто используемых датасетов для задач классификации болезней растений: он содержит десятки тысяч изображений листьев различных культур с метками здоровья/болезни и классами патологии. В оригинальной версии датасета содержится порядка $\approx$ 54.000 изображений, охватывающих примерно 38 классов заболеваний и здоровых состояний для 14 видов растений.

Изображения в датасете сделаны на нейтральном фоне и в стандартных условиях, что позволяет моделям более стабильно и эффективно обучаться, но ограничивает их способность к обобщению на реальные аграрные условия.

Датасет PlantVillage достаточно качественный и проверенный множеством исследований, однако содержит только метки классов, без области/bounding box'а заболевания, поэтому не подойдет для сегментации.

\subsection{PlantSeg ~\cite{plantseg}}

\textbf{PlantSeg} - это крупномасштабный датасет для задачи сегментации пораженных участков растений, созданный с учетом реальных ("in-the-wild") условий съемки. В отличие от большинства других датасетов, содержащих лишь метки классов или ограниченные bounding-box разметки, PlantSeg предоставляет детальные маски сегментации для каждого изображения, что делает его уникальным ресурсом для обучения и оценки моделей семантической сегментации.

Кроме качественной маски заболевания и диких условий у PlantSeg также:
\begin{itemize}
     \item \textit{Высокое качество изображения} - среднее 500x500+ разрешение изображения, против например 256x256 у PlantVillage
     \item \textit{Большее число классов и растений} - 115 класса и 34 вида растений против 38 и 14 у PlantVillage
\end{itemize}

PlantSeg очень хорошо подходит под нашу задачу и будет основным источником данных для обучения в дальнейших экспериментах.

\section{Цели и дальнейший план}

Дальнейшая работа будет состоять из двух основных частей - увеличение качества и метрик, и оптимизация данных для few-shot learning'а. Пока запланировано:
\begin{itemize}
\item Эксперименты с обучением DINOv3 ~\cite{dinov3}: подобрать разные архитектуры голов-декодеров (базовые сверточные, U‑Net‑подобные, т.д.), провести серию экспериментов с оптимизаторами, lr‑scheduler’ами и подобрать параметры обучения; обучить с большим числом эпох; оценить прирост по mIoU и mAP.

\item Метрики и анализ ошибок: расширить набор метрик качества для сегментации и классификации, провести визуализацию ошибок на валидационных примерах.

\item Попробовать также модель SAM (Segment Anything Model) ~\cite{sam3} для сегментации и сравнить с DINOv3, в сыром виде и с fine-tune.

\item Для оптимизации данных (а именно минимизации числа данных для обучения) изучить и построить бейзлайн с k-means / k-medoids по patch embeddings и Farthest Point Sampling, поискать дополнительные работы

\end{itemize}

% \subsection{Рисунки}

% \begin{figure}[ht]
% 	\centering
% 	\includegraphics[width=0.8\textwidth]{example.png}
% 	\caption{Пример графика. Тут должна быть подпись, поясняющая что происходит на рисунке (краткая, но достаточная для понимания основной идеи графика).}
% 	\label{fig:by_epochs}
% \end{figure}

% Все рисунки в тексте должны иметь подписи и вы на них должны ссылаться в тексте. Например, на Рисунке~\ref{fig:by_epochs} изображен пример графика. Не забывайте подписывать все оси на графиках, добавлять легенду и пояснять все обозначения, а также используйте адекватного размера шрифты и толщину линий на графиках (все должно быть видно и понятно без многократного увеличения). На рисунке из примера явно не хватает обозначения синей линии в легенде.


% \subsection{Таблицы}

% Все таблицы в тексте тоже должны иметь подписи и вы на них должны ссылаться в тексте. Например, в Таблице~\ref{table:long_epochs} показаны результаты примерного эксперимента. 


% \begin{table}[ht]
% 	\caption{Пример таблички. Тут должна быть подпись, поясняющая что происходит в таблице (краткая, но по делу).}
% 	\label{table:long_epochs}
% 	\footnotesize
% 	\centering
% 	\begin{tabular}{lrrrrrrrr}
% 		\toprule
% 		& \multicolumn{3}{c}{$\mathsf{Val}$} &
% 		\multicolumn{3}{c}{$\mathsf{Test}$} \\
% 		\cmidrule(lr){2-4} \cmidrule(l){5-7} 
% 		{} &  $\mathsf{Prec}$ &  $\mathsf{Rec}$ &  $\mathsf{F1}$ &  $\mathsf{Prec}$ &  $\mathsf{Rec}$ &  $\mathsf{F1}$  &  $\mathsf{nodes}$ & $\mathsf{subtokens}$\\
% 		\midrule
% 		запуск 1    &    0.4894 &   0.3775 &  0.4263 &     0.4824 &    0.3683 &   0.4177 & 10029 & 179\\
% 		запуск 2    &    0.4887 &   0.3739 &  0.4237 &     0.4891 &    0.3724 &   0.4228 & 10039 & 177\\
% 		запуск 3    &    0.4820 &   0.3751 &  0.4219 &     0.4838 &    0.3677 &   0.4178 & 10037&	180\\
% 		\midrule
% 		\bf{среднее} &    \bf{0.4867} &   \bf{0.3755} &  \bf{0.4239} &    \bf{ 0.4851} &    \bf{0.3695} &   \bf{0.4195} \\
% 		\bf{дисперсия}  &    0.0041 &   0.0019 &  0.0022 &     0.0036 &    0.0025 &   0.0029 \\
% 		\bottomrule
% 	\end{tabular}
% \end{table}

% \subsection{Формулы}

% Формулы стоит центрировать, а также нумеровать, если вы ссылаете на них в тексте. Также не забывайте пояснять все обозначения в формулах. Например, запишем следующую задачу оптимизации:
% \begin{equation}
%     \label{eq:si_opt}
%         \theta* = \min_{\theta} F(\theta),
% \end{equation}
% где $F$~-- квадратичная функция от параметра $\theta$. При необходимости, далее в тексте можно сослаться на формулу~(\ref{eq:si_opt}). При этом, в зависимости от конкретных формул, можно использовать разные слова: формула, уравнение, задача оптимизации и т.п.


	
\newpage 
\printbibliography[heading=bibintoc] 

% \begin{thebibliography}{0}
% 	\bibitem{chirkova18}\hypertarget{chirkova18}{}
% 	\href{https://arxiv.org/abs/1810.10927}
% 	{Nadezhda Chirkova, Ekaterina Lobacheva, Dmitry Vetrov. Bayesian Compression for Natural Language Processing. In EMNLP 2018.}
% \end{thebibliography}
	
% \newpage
% \appendix

% \section{Пример секции аппендикса}

\end{document}
